{\small\textbf{Why not the Basel rule? (Diego, 1:15) [8:55–10:10]}}\par\smallskip
Thank you. Let me start with the rule DNB chose \emph{not} to follow.
\begin{itemize}\setlength\itemsep{0pt}
\item \textcolor{navy}{\textbf{[def]}} The Basel rule looks at the credit-to-GDP gap: how far total credit, as a share of GDP, sits above its long-run trend. That trend is estimated with a statistical smoother called the Hodrick-Prescott filter.
\item A gap below two points means a zero buffer; above ten points means the maximum.
\end{itemize}
\par\smallskip
We rebuilt this indicator ourselves and match the official BIS series almost exactly. On Dutch data the rule has a poor record:
\begin{itemize}\setlength\itemsep{0pt}
\item \textbf{It missed the last crisis:} the gap was minus 13 points on the eve of the financial crisis.
\item \textbf{It raised a false alarm} in 2012, during a house-price bust, simply because GDP was falling.
\item \textbf{It gets rewritten after the fact:} in 2016 it read slightly negative in real time, but plus 27 points with today's data.
\end{itemize}
\par\smallskip
Why does it fail so badly in the Netherlands? Dutch credit rose from under fifty percent of GDP in the 1960s to around 350 percent by 2015. The filter interprets much of that structural deepening as "trend". When the ratio later comes down, as it has since its peak of about 350 percent in 2015, the gap turns deeply negative, even if nothing about cyclical risk has improved.
\par\smallskip
At the two decisions it read minus 33 and minus 47 points, so the rule said zero. On this record, DNB had good reason to use a different compass.
